\documentclass[11pt, twocolumn]{article}
\author{Sach Shrestha \\
        \small University of Georgia}
\date{October 5, 2025}
\title{\textbf{Antibiotic Resistance Report}}
\usepackage{graphicx}
\usepackage[labelfont=bf, font=small]{caption}
\usepackage{natbib}
\usepackage[T1]{fontenc}
\usepackage{lmodern}



\providecommand{\keywords}[1]
{
  \small	
  \textbf{\textit{Keywords---}} #1
}


\begin{document}

\twocolumn[
	\begin{@twocolumnfalse}
	\maketitle
	\centering
		\begin{abstract}
    		Antibiotics have been used to treat bacterial infections and fight bacteria since the beginning of the 20th century. It has become an indispensable tool that has saved countless lives. However, there is recent evidence that suggests bacteria are developing resistance against these antibiotics through environmental overexposure caused by pollution and misuse of drugs. If these bacteria continue developing stronger resistances, it will render our current antibiotics deficient in their treatment. Thus, the goal of this BIOL1103L study is to examine the performance of our current line of antibiotic/ antiseptic drugs against two specific strains of bacteria: \textit{Escherichia coli} (\textit{E. coli}) and \textit{Bacillus subtilis (B. subtilis)}. It was hypothesized that, in general, gram positive antibiotics would be effective against\textit{ B. subtilis}, which is a gram-positive bacterium; and gram-negative antibiotics would be effective against \textit{E. coli}, which is gram negative. The way we measured effectiveness of each drug was by measuring the zone-of-inhibition (ZOI) expressed by each when placed on a bacterial lawn containing either \textit{E. coli} or \textit{B. subtilis} sample. Our hypothesis was largely found to be accurate; but a slightly alarming discovery, however, was that \textit{E.coli} resisted majority of the drugs we tested.  This could be because of our small sample size for antibiotics and number of test cases per antibiotic, which might not be able to accurately represent the entire population. We suggest more extensive research be conducted on the subject and bacterial resistances be monitored so we do not fall behind in providing effective treatment.
    		\end{abstract} \hspace{10pt}
    		
    	
%TC:ignore
\keywords{antibiotic resistance, zone-of-inhibition, bacteria lawn, \textit{E. coli}, \textit{B. subtilis}}    	
%TC:endignore
    	
	\end{@twocolumnfalse}
]

\clearpage

\section{Introduction}

Antibiotics are widely used drugs that help treat bacterial infections. The rampant usage of these drugs has made their occurrence ubiquitous in the environment \citep{Gothwal_2015} . According to \citet{Akova_2016}, broad spectrum antibiotics in particular are liberally and most unnecessarily used, and such practices have caused dramatic increases in emerging resistances (p. 252). This emerging resistance among bacteria to commonly used antibiotics and last resort drugs (such as vancomycin) is the most concerning problem \citep{Urban-Chmiel_2022}. It is a serious global health concern that persists despite numerous and varied efforts to improve circumstances \citep{Velez_2016}. 

As such, we are interested in examining the current level of resistances possessed by two specific strains of bacteria:  \textit{Escherichia coli} (\textit{E. coli}) and \textit{Bacillus subtilis (B. subtilis)}, against a roster of some commonly used antibiotics and antiseptics. We hypothesize that antibiotics will be most effective against bacteria of the same gram strain (positive or negative). We will measure effectiveness of each drug by measuring the zone-of-inhibition (ZOI) expressed by each when placed on a bacterial lawn containing either \textit{E. coli} or \textit{B. subtilis} sample. For the purposes of this study, we will consider drugs with ZOI $<$ 14mm to be ineffective in their treatment, that is to say the bacteria is resistant to the drug.

\section{Method}

We performed plate streaks of \textit{E. coli} and \textit{B. subtilis} bacteria onto their respective agar plates using cotton swabs, creating a bacteria lawn. This lawn was then further divided into 4 parts, a quarter zone for each drug. Next, we placed paper discs dipped in antibiotic/ antiseptic solution into their respective zones in the plate. A disc dipped in distilled water was placed in the center to serve as our control sample. Then, we sealed the plates and put them aside to let the bacteria in the plates grow. A week later, we measured the zone of inhibition (ZOI) expressed by each disc vs. the bacteria growth and tallied up the results in a database for analysis. Finally, we calculated the average ZOI (mm) for each drug to get the most accurate idea of its effectiveness.

\section{Results}
\subsection{\textit{E. coli}}

\begin{figure}[h]
	\includegraphics[width=\columnwidth]{figure1.png} 
	\caption{The Zone of Inhibition ZOI (mm) for each drug vs. \textit{E. coli} bacteria. Pairs of students perform the experiment by testing 5 Antibiotics/ Antiseptics of their choosing, the measurements for the ZOI of each drug after 1 week of bacterial growth are then tallied in a database and averaged out across the entire class to get the final ZOI results presented in this graph. Larger ZOI suggests the drug is more effective  against\textit{E. coli}. Smaller ZOI suggests that \textit{E. coli} is more resistant to the drug.}
	\label{fig:ecol1}
\end{figure}

\begin{figure}[h]
	\includegraphics[width=\columnwidth]{figure3.png}  
	\caption{Drug type ZOI (mm) vs. \textit{E. coli} bacteria. The colored lines represent the mean ZOI for each drug type: Antiseptic (red), Gram - (green), and Gram + (blue). The black line represents the antibiotic-resistant threshold of ZOI $<$ 14mm. A mean ZOI below the black line would suggest \textit{E. coli} possesses some resistance against that category of drug.}
	\label{fig:ecol2}
\end{figure}

The mean ZOI (mm) for each drug vs. \textit{E. coli} is displayed in Figure~\ref{fig:ecol1}. The antibiotics tested against it were as follows: Amoxicillin-clavulanic acid (broad gram +), Cefazolin (broad gram +), Gentamicin (narrow gram -), Penicillin (broad gram +), Tetracycline (broad gram +/-, for this study we consider it gram -), and Vancomycin (narrow gram +). The antiseptics were as follows: Chlorox, Mrs. Meyer’s, Listerine, and Pine Sol. We found that the mean average ZOI for E. coli treated with antibiotics was 80\% larger than the mean average ZOI for E. coli treated with antiseptics. Furthermore, using the interpretation that ZOI $<$ 14mm = resistant, we noted that E. coli was resistant to 70\% of the antibiotics/ antiseptics tested, which includes 50\% of the antibiotics and every single antiseptic. Among the antibiotics that were effective against E. coli, the average ZOI produced by the most effective drug, Gentamicin, was found to be 38\% larger than the next best performing drug, Cefazolin. Finally, Mrs. Meyer’s solution was the least effective drug against E. coli, with it producing no ZOI, and least effective antibiotic was Vancomycin.

In Figure~\ref{fig:ecol2}, we separated antibiotics further by their respective gram strains, gram + and gram -, to examine their effectiveness against \textit{E. coli}. We found that on average, gram - antibiotics performed 38\% better than gram + antibiotics. We can also see that on average, both antiseptics and gram + antibiotics appear to be resisted by the \textit{E. coli} bacteria.


\subsection{\textit{B. subtilis}}

\begin{figure}[h]
	\includegraphics[width=\columnwidth]{figure2.png} 
	\caption{The Zone of Inhibition ZOI (mm) for each drug vs. \textit{B. subtilis} bacteria. Pairs of students perform the experiment by testing 5 Antibiotics/ Antiseptics of their choosing, the measurements for the ZOI of each drug after 1 week of bacterial growth are then tallied in a database and averaged out across the entire class to get the final ZOI results presented in this graph. Larger ZOI suggests the drug is more effective against \textit{B. subtilis}. Smaller ZOI suggests that \textit{B. subtilis} is more resistant to the drug.}
	\label{fig:bsub1}
\end{figure}

\begin{figure}[h]
	\includegraphics[width=\columnwidth]{figure4.png} 
	\caption{Drug type ZOI (mm) vs. \textit{B. subtilis} bacteria. The colored lines represent the mean ZOI for each drug type: Antiseptic (red), Gram - (green), and Gram + (blue). The black line represents the antibiotic-resistant threshold of ZOI $<$ 14mm. A mean ZOI below the black line would suggest \textit{B. subtilis} possesses some resistance against that category of drug.}
	\label{fig:bsub2}
\end{figure}

The mean ZOI (mm) for each drug vs. \textit{B. subtilis} is displayed in Figure~\ref{fig:bsub1}. The antibiotics tested against it were as follows: Amoxicillin-clavulanic acid (broad gram +), Cefazolin (broad gram +), Gentamicin (narrow gram -), Penicillin (broad gram +), Tetracycline (broad gram +/-, for this study we consider it gram -), and Vancomycin (narrow gram +). The antiseptics were as follows: Chlorox, Mrs. Meyer’s, Listerine, Pine Sol, and Tom’s of Maine. We found that the mean average ZOI for B. subtilis treated with antibiotics was 72\% larger than the mean average ZOI for B . subtilis treated with antiseptics. Furthermore, using the interpretation that ZOI $<$ 14mm = resistant, we noted that B. subtilis was susceptible to over 90\% of the antibiotics/ antiseptics tested. Among the drugs that were effective against B. subtilis, the average ZOI produced by the most effective drug, Cefazolin, was found to be 40\% larger than the next best performing drug, Gentamicin. Finally, Tom’s of Maine was the least effective drug against B. subtilis, with it producing no ZOI.

In Figure~\ref{fig:bsub2}, we separated antibiotics further by their respective gram strains, gram + and gram -, to examine their effectiveness against \textit{B. subtilis}. We found that on average, gram + antibiotics performed 15\% better than gram - antibiotics. We can also see that on average, all drug types effectively inhibited \textit{B. subtilis} growth to at least a moderate extent. Furthermore, cross referencing their performances; we noted that on average, gram + antibiotics performed 2.3x better against \textit{B. subtilis} than they did against \textit{E. coli}, and gram - antibiotics performed 46\% better against \textit{B. subtilis} than they did against \textit{E. coli}.


\section{Discussion}

As noted in results section, we saw antibiotics of each gram strain perform better against their respective gram strain bacteria. This was expected and follows our hypothesis. Interestingly, Gentamicin, the narrow gram – antibiotic performed well against the gram – bacteria \textit{E. coli}, but it also excelled against the gram + bacteria \textit{B.subtilis}. \textit{E. coli}  proved to be resistant to 70\% of drugs tested, which includes all of the antiseptics tested. It can be concluded that the sampled strain of \textit{E. coli} is very resilient and advanced in its resistances. This might suggest that \textit{E. coli} could be on its way to further develop resistances to the remaining drugs it is currently susceptible to. 

\textit{B. subtilis} on the other hand, appears to be weak to 91\% of the drugs tested. It was found to even be susceptible to gram – antibiotics and antiseptics. This could suggest that the sampled strain of \textit{B. subtilis} is not as resilient or advanced as the \textit{E. coli} sample.
  
We can use this information to formulate course of action to fight antibiotic resistance in both \textit{E. coli} and \textit{B. subtilis}. For strains as advanced and resilient as \textit{E. coli}, as \citet{Urban-Chmiel_2022} suggest, there is an urgent necessity to implement and develop methods of bacterial control alternative to antibiotics, along with research into prevention of further resistance development. For strains still developing resistances like the \textit{B. subtilis} sample in this study, it is essential to invest in more research to prevent these resistances from building.  Since the problem of emerging resistance is multifaceted, its prevention and control will require multiple, coordinated interventions by various parties \citep{Akova_2016} (p. 259), so it is vital to raise awareness among the public and those involved.  


\section{Conclusion}

Antibiotic resistance is a disturbing phenomenon with deadly consequences. As such it is crucial to monitor resistance levels among common bacterial strains. We examined a common line-up of bacterial control drugs and as hypothesized, the best treatment for \textit{E. coli} and \textit{B. subtilis} are antibiotics of matching gram strains. \textit{E. coli} was more resilient to treatment than \textit{B. subtilis}, with it resisting majority of the drugs tested including broad-spectrum antibiotics and antiseptics when \textit{B. subtilis} was found to be susceptible to almost all drugs tested. 

While these results help paint a picture of the current levels of antibiotic resistance in \textit{E. coli} and\textit{ B. subtilis}, the clear limitations of the study mean the results may not accurately represent the entire population of gram +, gram -, and broad-spectrum antibiotics; the sample size of drugs and number of test cases per drug we tested were on the smaller side, and the time frame of a week of bacterial growth may not be suitable to gauge effectiveness of each drug. 


\bibliography{references}
\bibliographystyle{apalike}



\end{document}